\documentclass{beamer}
\usepackage{beamerthemeOerestad}
\usetheme{Oerestad}

\usepackage[utf8]{inputenc}
\usepackage{amsfonts}
\usepackage{amsmath}
\usepackage{booktabs}
\usepackage{color}
\usepackage{colortbl}
\usepackage{graphicx}
\usepackage{hyperref}
\usepackage{listings}
\usepackage[square,sort,comma,numbers]{natbib}
\usepackage{url}
\usepackage{semantic}
\usepackage{tikz}
\usepackage{todonotes}

\lstset{
  language=Lisp,
  basicstyle=\ttfamily,
  otherkeywords={define,match,let,:,->},
  identifierstyle=\color{pink!60!blue},
  keywordstyle=\color{blue!50!black},
  stringstyle=\color{green!50!black},
  commentstyle=\color{pink!60!black}
}

\title{Functional Programming in Typed Racket}
\author{Florian Florian \\\small{\texttt{fbie@itu.dk}} \\~}

\institute{IT University of Copenhagen \& UCAS}

\date{2016-05-26}

\begin{document}

\begin{frame}
  \titlepage{}
\end{frame}

\begin{frame}{Who am I?}
  Just call me Florian!

  \begin{itemize}
  \pause{} \item I am from Germany and live in Denmark.
  \pause{} \item Second year PhD student at IT University of Copenhagen and UCAS.
  \pause{} \item My Chinese is really bad (but I am trying to learn!)
  \pause{} \item If you do not understand what I say or if I speak too fast, \textbf{please tell me right away!}
  \end{itemize}
\end{frame}

\begin{frame}{Why Functional Programming?}

  Functional programming is old but recently it has become very popular!

  \begin{itemize}
  \pause{} \item Makes it easier to think about what a program does.
  \pause{} \item Good for parallel programming, performance (see next lecture).
  \pause{} \item Used in financial industry, modern compilers and more.
  \pause{} \item Jobs in functional programming pay better!
  \end{itemize}

\end{frame}

\begin{frame}{Typed Racket}

  \begin{center}
    \includegraphics[width=0.2\textwidth]{racket.png}
  \end{center}

  \textbf{Racket} and its sister-language \textbf{Typed Racket} are Scheme-dialects.

  \begin{itemize}
  \pause{} \item As the names say, Racket is untyped, Typed Racket has types.
  \pause{} \item Functional language but with support for objects.
  \pause{} \item Great meta-programming (we won't look at that).
  \end{itemize}
\end{frame}

\begin{frame}[fragile]{Hello World in Typed Racket}
\begin{center}
\begin{lstlisting}
;; Tell run-time which language to use.
#lang typed/racket

;; Now print something.
(print "Nihao!")
\end{lstlisting}
\end{center}
\end{frame}

\begin{frame}[fragile]{A More Interesting Program}
\begin{center}
\begin{lstlisting}
#lang typed/racket

;; Define x as the result of the computation.
(define x (* (+ 2 4) (+ 42 9)))
\end{lstlisting}
\end{center}
\end{frame}

\begin{frame}[fragile]{Do You Think This Looks Strange?}

  In Racket, all expressions are written like this:

  \begin{center}
    \lstinline{(function arg1 arg2 ... argN)}
  \end{center}

  \pause{}

  Operators are also just functions:

  \begin{center}
    \begin{tabular}{ccc}
      \lstinline{(+ x y)} & $\Rightarrow$ & $x + y$ \\
      \lstinline{(> x y)} & $\Rightarrow$ & $x > y$ \\
      \lstinline{(/ x y)} & $\Rightarrow$ & $\frac{x}{y}$ \\
      \lstinline{(f x y)} & $\Rightarrow$ & $f(x, \, y)$
    \end{tabular}
  \end{center}

  \pause{}

  Some functions are \textit{variadic}:

  \begin{center}
    \begin{tabular}{ccc}
      \lstinline{(+ x y z)} & $\Rightarrow$ & $x + y + z$
    \end{tabular}
  \end{center}

\end{frame}

\begin{frame}[fragile]{Local Bindings}

Just like local variables in Java, but you can never change them!

\begin{lstlisting}
#lang typed/racket

(let ([x (* 2 16)]
      [y (* 3 17)])
    (+ x y))
\end{lstlisting}

What does this program do?

\vspace{1cm}

\textit{Note:} You cannot reference \lstinline{x} and \lstinline{y} after the last closing parenthesis of the \lstinline{let} expression!

\end{frame}

\begin{frame}[fragile]{A First Function}
\begin{lstlisting}
#lang typed/racket

(: times-two (-> Number Number))
(define (times-two x)
  (* 2 x))
\end{lstlisting}

\pause{}

Several new things on this slide:

\begin{itemize}
\item Functions need no return statement. Their return value is the last executed statement!
\item Type annotations start with \lstinline{:} and describe the type of a symbol.
\item The type of \lstinline{times-two} is $Number \, \rightarrow \, Number$.
\end{itemize}
\end{frame}

\end{document}
